
Cross-architecture property prediction from neural network weights enables automated model selection and architecture search without costly evaluation. This work addresses the performance gap between within-architecture prediction ($\rho$ = 0.81) and cross-architecture prediction ($\rho$ = 0.54) by decomposing architectural differences into three learnable signals. CAPE (Cross-Architecture Parameterized Encoder) employs operation-specific encoders for convolutional, attention, and MLP operations; contrastive projection to align heterogeneous embeddings via task objectives; and architecture graph neural networks to capture computational topology. In proof-of-concept validation with 100 models (50 ResNet-50, 50 ViT-Base) and synthetic ImageNet accuracy labels, CAPE achieved $\rho$ = 0.67 on $ResNet\rightarrow ViT$ transfer, a 24\% improvement over the Set Neural Encoder baseline ($\rho$ = 0.54, $\Delta\rho$ = 0.13, p = 0.032 via permutation test with 1000 iterations). Ablation studies show independent contributions from each component: operation encoders (+0.04), contrastive projection (+0.05), and GNN residual (+0.04). Diagnostic metrics confirm distinct operation representations (cosine similarity = 0.12 < 0.95), preserved architectural structure (intra-architecture variance = 0.15 $\geq$ 0.1), and meaningful topology contribution (learned residual weight $\alpha$ = 0.42 > 0.1). These results establish that architectural differences, when decomposed and aligned through shared objectives, enable mechanism validation for cross-architecture weight-space learning. Full-scale validation with 400 models across four architectures and real accuracy labels is deferred to future work.
